\section{Accuracy and state size across descriptions of training}
\label{sec:resources}

A common resource question connects otherwise different theories: how much
evolving information suffices to approximate the same learning process?
The first distinction is between simplifying that process and compressing its
description.  Linear, shallow and frozen-feature models give tractable accounts
of important regimes.  Dense training, tangent hierarchies, limiting-response
solvers and response memory instead provide different representations to compare
when depth, nonlinearity and feature learning are retained
(\Cref{tab:resource-comparison}).

\paragraph{A shared target.}
Fix the architecture, data, initialization law (including the readout), canonical
scaling and a population flow $f_\infty$ on $[0,T]$.  For a test probability
measure $Q$, write
\begin{equation}
 E_Q(\widehat f;T)=\sup_{t\le T}
 \norm{\widehat f(t,\cdot)-f_\infty(t,\cdot)}_{L^2(Q)},\qquad
 \mathcal E_Q=(\E E_Q^2)^{1/2}.
 \label{eq:resource-error}
\end{equation}
The expectation covers initialization or numerical sampling.  Uniform angular
$Q$ gives whole-circle prediction RMSE.  The reverse triangle inequality bounds
the difference between the two predictors' test RMSEs by $E_Q$.
This common observable allows comparison with methods that predict the learned
function without reconstructing all trained weights; our parameter theorem
also supplies the latter.

\paragraph{Matched-regime idealizations.}
The table separates representation counts from conditional population accuracy.
For the quantitative comparison, assume the specified population target exists,
the dense error is bounded by $A_T/\sqrt n$, and the closure estimates are
uniform in width in a normalized prediction-controlling norm.  Then coupling
the closure and dense network through the same initialization gives
\begin{equation}
 \mathcal E_{\rm dense}\le A_Tn^{-1/2},\qquad
 \mathcal E_q\le A_Tn^{-1/2}+B_{q,T}P^{-q},\quad q=1,2.
 \label{eq:resource-laws}
\end{equation}
Here $q=1$ denotes the learning-speed clock and $q=2$ the joint clock.
The root-width scale is supported by perturbative fluctuation theory
\citep{bordelon2023finite}; a general matched rate and width-uniform closure
constants are additional hypotheses, not conclusions of
\Cref{thm:old,thm:new}.  All constants can depend on data, $m,H,T$ and $Q$.
Numerical limiting-process bounds additionally require stable sampling,
time discretization and self-consistency.  These assumptions are explicit
because they determine what can be compared, rather than supplying a universal
ordering of methods.

\begin{table}[!t]
\centering
\small
\setlength{\tabcolsep}{4pt}
\setlength{\belowcaptionskip}{6pt}
\renewcommand{\arraystretch}{1.18}
\caption{Descriptions of training on a common state--accuracy axis.
The upper block changes a property of the target dynamics; its rows are not
general accuracy refinements of the specified deep feature-learning flow.
The lower block gives direct representation counts and conditional
population-prediction laws explained in the text.  Moving state excludes
fixed resources consistently.  Fixed input dimension and lower-order outer
weights are suppressed; $m\ge2$ for the direct NTH count.}
\label{tab:resource-comparison}
\begin{tabular}{@{}>{\raggedright\arraybackslash}p{0.225\linewidth}>{\raggedright\arraybackslash}p{0.205\linewidth}>{\raggedright\arraybackslash}p{0.22\linewidth}>{\raggedright\arraybackslash}p{0.285\linewidth}@{}}
\toprule
Description & Moving state & Fixed resources & Error mechanism or regime change\\
\midrule
\multicolumn{4}{@{}l}{\emph{Simplifying the training regime}}\\
Deep linear & --- & --- & Removes nonlinear activations\\
One-hidden-layer mean field & --- & --- & Removes multilayer hidden interactions\\
Frozen NTK & $m$ & $m^2$ kernel; test evaluator &
Feature-freezing bias relative to rich training\\
\midrule
\multicolumn{4}{@{}l}{\emph{Representations of deep nonlinear feature learning}}\\
Dense width $n$ & $O(Hn^2)$ & --- & $A_Tn^{-1/2}$\\
NTH through $p$ & $\sum_{j=1}^{p-1}m^j$ &
$m^p$ top tensor; mixed-kernel functions &
$\eta_p(T)$ from higher-kernel bounds\\
Sampled TP/DMFT, full grid & $O(Hm^2K^2)$ &
Model and input covariances &
$A_{{\rm D},T}S^{-1/2}$\newline
$+B_{{\rm D},T}(T/K)^r+\eta$\\
Learning-speed memory & $O(HmnP)$ &
$O(Hn^2)$ initialization &
$A_Tn^{-1/2}+B_{1,T}/P$\\
Joint response memory & $O(HmnP+P^2)$ &
$O(Hn^2)$ initialization; prefixes &
$A_Tn^{-1/2}+B_{2,T}/P^2$\\
\bottomrule
\end{tabular}
\end{table}

\paragraph{What each refinement buys.}
NTH evolves features through successively higher tangent kernels.  Its
same-target truncation error $\eta_p(T)$ depends on their growth; the published
width-dependent theorem cannot simply be transferred to another scaling.
\Cref{app:comparison-nth} derives a feedback bound for the canonical hierarchy
and explains why it is not a time-Taylor truncation.
For TP/DMFT we specify a sampled full-grid response solver
\citep{yang2021tensor,bordelon2022dmft}: $K$ is the number of time steps,
$S$ the number of paths, $r$ the numerical time order and $\eta$ the solver
error.  Streaming paths leaves $O(Hm^2K^2)$ kernel state while increasing
sampling work.  Thus $S$ is not a finite-network width, and a numerical
time grid is not the physical horizon $T$.

\paragraph{Size at prescribed accuracy.}
Under \eqref{eq:resource-laws}, optimizing width and order at fixed task gives
the following leading moving-memory scalings:
\begin{align}
 M_{\rm dense}^*(\varepsilon)&\asymp
 HA_T^4\varepsilon^{-4},\nonumber\\
 M_{\rm old}^*(\varepsilon)&\asymp
 HmA_T^2B_{1,T}\varepsilon^{-3},\qquad
 M_{\rm joint}^*(\varepsilon)\asymp
 HmA_T^2\sqrt{B_{2,T}}\,\varepsilon^{-5/2}.
 \label{eq:resource-optima}
\end{align}
These optimize sufficient error bounds for the stated representations;
\Cref{app:comparison-opt} gives the allocations and the corresponding
conditional NTH and numerical TP/DMFT costs.  They are not lower bounds on
all algorithms within a framework.  Keeping dense initialization leaves a
leading $\varepsilon^{-4}$ contribution to our total storage and arithmetic.
The established gain is in the evolving learned state; improving total cost
also requires treating the fixed environment.

This comparison separates width resolution from history resolution.
At fixed width the response-memory rates are proved, and the state is linear
in the number of retained moments.  The population calculation predicts how
that gain combines with width approximation.  It does not assume equal or
small horizon constants across methods: one fixed order serves every time
inside a chosen interval, but increasing that interval may require more
moments.  \Cref{app:comparison} records the derivations, per-step costs and
test-function evaluation requirements so that these common-axis comparisons
can be refined without changing the underlying question.
